%! Inverse Trigonometric Functions — Unit 7, Lesson 7.2 — Lesson Plan
\documentclass[10pt]{article}
\usepackage{precalculus-article}
\usepackage{precalculus-boxes}
\usepackage{graphicx}
\graphicspath{{images/}}
\newcommand{\TallMath}[1]{$\displaystyle #1\rule[-1.4em]{0pt}{3.2em}$}
\newcommand{\CourseName}{Precalculus}
\newcommand{\SchoolYear}{2026--2027}
\newcommand{\MeetingLength}{55 minutes}
\newcommand{\UnitNumberName}{Unit 7: Analytic Trigonometry and Polar Coordinates \quad}
\newcommand{\LessonNumberName}{Lesson 7.2: Inverse Trigonometric Functions}

\begin{document}
\begin{center}
    {\Large\bfseries \CourseName: \SchoolYear \\
    {\small\bfseries \UnitNumberName \LessonNumberName}}
\end{center}

\begin{tcolorbox}[colback=sky, colframe=navy, boxrule=0.9pt, arc=2mm,
    left=2mm, right=2mm, top=1.5mm, bottom=1.5mm]
    \textbf{Primary Objective:} Students construct analytical and graphical
    representations of the inverse sine, cosine, and tangent functions over
    restricted domains, and interpret output values as angle measures.
    (2--3 instructional periods, \MeetingLength\ each)
\end{tcolorbox}

\renewcommand{\arraystretch}{1.6}
\begin{skillbox}[Priority Ideas \& Skills]{goldbox}
\begin{minipage}[t]{0.54\linewidth}\vspace{0pt}
    \textbf{AP Skills}
    \begin{itemize}[leftmargin=1.4em, itemsep=2pt]
        \item \textbf{1.C} --- Construct new functions (inverses)
        \item \textbf{2.B} --- Construct equivalent representations
    \end{itemize}
    \textbf{Learning Objective}
    \begin{itemize}[leftmargin=1.4em, itemsep=2pt]
        \item \textbf{3.9.A} --- Construct analytical and graphical
              representations of inverse sine, cosine, and tangent over
              restricted domains
    \end{itemize}
\end{minipage}\hfill
\begin{minipage}[t]{0.42\linewidth}\vspace{0pt}
    \textbf{Key Understandings (EKs)}
    \begin{itemize}[leftmargin=1.4em, itemsep=2pt]
        \item Input/output switch: output of inverse trig $=$ angle measure
        \item Trig functions are periodic (many-to-one) $\Rightarrow$ must
              restrict domain to invert
        \item Restricted domains: $\sin\colon[-\tfrac{\pi}{2},\tfrac{\pi}{2}]$,
              $\cos\colon[0,\pi]$, $\tan\colon(-\tfrac{\pi}{2},\tfrac{\pi}{2})$
    \end{itemize}
\end{minipage}
\end{skillbox}

\begin{skillbox}[{Vocabulary, Concepts, \& Theorems}]{greenbox}
\begin{multicols}{2}
\begin{itemize}[leftmargin=1.4em, itemsep=2pt]
    \item \textbf{Arcsine} ($\arcsin x$ or $\sin^{-1}x$) --- inverse of the
          restricted sine function
    \item \textbf{Arccosine} ($\arccos x$ or $\cos^{-1}x$) --- inverse of the
          restricted cosine function
    \item \textbf{Arctangent} ($\arctan x$ or $\tan^{-1}x$) --- inverse of the
          restricted tangent function
    \item \textbf{Restricted domain} --- a subset of the natural domain chosen
          so the function is one-to-one
    \item \textbf{Principal value} --- the unique output angle in the
          restricted range of an inverse trig function
    \item \textbf{One-to-one function} --- each output value corresponds to
          exactly one input value (passes horizontal line test)
\end{itemize}
\end{multicols}
\end{skillbox}

\begin{skillbox}[Activate Prior Knowledge \& Spiral Review]{sky}
\begin{tabularx}{\linewidth}{@{}X c@{}}
    \begin{minipage}[t]{0.60\linewidth}\vspace{0pt}
        \textbf{Warm-Up reviews:}
        \begin{itemize}[leftmargin=1.4em, itemsep=2pt]
            \item Evaluating $\sin\theta$, $\cos\theta$, $\tan\theta$ at
                  multiples of $\tfrac{\pi}{6}$ and $\tfrac{\pi}{4}$
            \item Domain and range vocabulary for a function
            \item Identifying one-to-one functions (horizontal line test)
                  from a table of values
        \end{itemize}
    \end{minipage}
    &
    \begin{minipage}[t]{0.34\linewidth}\vspace{0pt}\centering
        See attached warm-up.
    \end{minipage}
\end{tabularx}
\end{skillbox}

\begin{skillbox}[Hook]{sky}
    Display: \textit{``A ramp rises 3 feet over a horizontal run of 5 feet.
    What angle does the ramp make with the ground?''}

    Students quickly realize they can compute $\tan\theta = \tfrac{3}{5}$, but
    need to \emph{undo} the tangent to find $\theta$.  Ask: \textit{``What
    operation undoes addition?  What operation undoes tangent?''}  Introduce
    the idea of an inverse function for trigonometry.
\end{skillbox}

\begin{skillbox}[Lesson]{sky}
\begin{multicols}{2}
\textbf{Day 1 --- Why Restricted Domains?}
\begin{itemize}[leftmargin=1.4em, itemsep=2pt]
    \item Use the table of $\sin$ values to show that $\sin\tfrac{\pi}{6} =
          \sin\tfrac{5\pi}{6} = \tfrac{1}{2}$. Ask: can we define an inverse
          if two inputs give the same output?
    \item Introduce the horizontal line test; connect to one-to-one functions
          from Unit 1.
    \item Restrict: sine to $[-\tfrac{\pi}{2}, \tfrac{\pi}{2}]$, cosine to
          $[0,\pi]$, tangent to $(-\tfrac{\pi}{2}, \tfrac{\pi}{2})$.  Ask
          students why these particular intervals.
\end{itemize}

\columnbreak

\textbf{Day 2 --- Defining \& Evaluating Inverse Trig}
\begin{itemize}[leftmargin=1.4em, itemsep=2pt]
    \item Fill in the restricted-domain table in the notes (Section 2).
    \item Model evaluation: $\arcsin\tfrac{1}{2} = ?$ --- ``Which angle in
          $[-\tfrac{\pi}{2}, \tfrac{\pi}{2}]$ has sine equal to $\tfrac{1}{2}$?''
          Answer: $\tfrac{\pi}{6}$.
    \item Work through the notes evaluation table (Section 3) together.
    \item Discuss notation: $\sin^{-1}$ means \emph{inverse function}, not
          $\frac{1}{\sin}$.
\end{itemize}
\end{multicols}
\textbf{Day 3 --- Compositions}
\begin{itemize}[leftmargin=1.4em, itemsep=2pt]
    \item When does $\sin(\arcsin x) = x$?  (Always, for $x \in [-1,1]$.)
    \item When does $\arcsin(\sin x) = x$?  (Only for $x \in [-\tfrac{\pi}{2}, \tfrac{\pi}{2}]$.)
    \item Model a composite like $\cos(\arcsin\tfrac{3}{5})$ using a
          right-triangle diagram drawn on the board.
\end{itemize}
\end{skillbox}

\begin{skillbox}[Active Monitoring]{redbox}
\textbf{Circulate and check:}
\begin{itemize}[leftmargin=1.4em, itemsep=2pt]
    \item Notes Section 2: Are students writing the correct interval notation
          (brackets vs.\ parentheses for tangent)?
    \item Notes Section 3: Are students identifying the \emph{principal value}
          (angle in the restricted range), not just any angle?
    \item Cold-call prompts: ``Why can't $\arcsin$ output $\tfrac{5\pi}{6}$?''
          and ``What is the domain of $\arccos$?''
\end{itemize}
\end{skillbox}

\begin{skillbox}[Group Work \& Differentiation]{redbox}
\begin{itemize}[leftmargin=1.4em, itemsep=2pt]
    \item \textbf{Tier R (Remediate):} Use a provided table matching trig
          values to angles; evaluate $\arcsin$, $\arccos$, $\arctan$ by
          look-up.  Focus on reading the table correctly.
    \item \textbf{Tier A (Apply):} Determine domain and range of each inverse
          trig function; evaluate expressions like $\sin(\arccos\tfrac{3}{5})$
          using right-triangle reasoning.
    \item \textbf{Tier E (Extend):} Compose trig and inverse trig functions
          analytically; determine precisely when $\sin(\arcsin x) = x$ vs.\
          $\arcsin(\sin x) = x$ and explain in terms of restricted domains.
\end{itemize}
\end{skillbox}

\begin{skillbox}[Individual Work \& Assessment]{redbox}
\textbf{Exit Ticket (2 items, 1 page):}
\begin{enumerate}[leftmargin=1.6em, itemsep=2pt]
    \item Evaluate two inverse trig expressions (principal values).
    \item Explain in 1--2 sentences why domain restriction is necessary for
          defining inverse trig functions.
\end{enumerate}
\textbf{Assessment note:} Look for correct identification of the principal
value and clear reasoning about one-to-one behavior.
\end{skillbox}

\begin{skillbox}[Reinforcement \& Extension]{goldbox}
\begin{itemize}[leftmargin=1.4em, itemsep=2pt]
    \item \textbf{Homework:} 5--6 evaluation/domain-range problems,
          2 AP-style MC items, 1 composition extension problem.
    \item \textbf{Extension:} Prove that
          $\arcsin x + \arccos x = \tfrac{\pi}{2}$ for all $x \in [-1,1]$.
    \item \textbf{Preview Lesson 3.10:} Ask students to think about what it
          means to \emph{solve} $\sin x = \tfrac{1}{2}$ --- how many solutions
          are there?  We will use inverse trig as a tool for solving trig
          equations.
\end{itemize}
\end{skillbox}
\end{document}
